\section{Omezení lineární regrese}\label{sec:ai-omezeni-regrese}

Lineární regrese je jednoduchá, názorná a~dobře interpretovatelná. Právě její jednoduchost je však zároveň omezením. Před použitím je vhodné prohlédnout bodový graf i~rezidua a~ověřit, zda lineární popis odpovídá povaze dat.

\begin{figure}[ht]
    \centering
    \begin{tikzpicture}[every node/.style={font=\normalsize}]
    \begin{scope}
        \draw[diredge] (0,0)--(4.4,0);
        \draw[diredge] (0,0)--(0,3.5);
        \draw[red!75!black,very thick] (0.5,0.7)--(4,2.8);
        \foreach \x/\y in {0.5/0.5,1/0.7,1.5/1.2,2/1.8,2.5/2.6,3/3.1,3.5/3.3,4/3.2}
            \fill[blue!75!black](\x,\y) circle(2pt);
        \node[below] at (2.2,-0.45){nelinearita};
    \end{scope}
    \begin{scope}[xshift=5.8cm]
        \draw[diredge] (0,0)--(4.4,0);
        \draw[diredge] (0,0)--(0,3.5);
        \draw[red!75!black,very thick] (0.5,0.5)--(4,3.1);
        \foreach \x/\y in {0.6/0.7,1.1/1,1.7/1.3,2.3/1.6,2.9/1.9,3.5/2.2}
            \fill[blue!75!black](\x,\y) circle(2pt);
        \fill[orange!85!black](3.9,0.4) circle(3pt);
        \node[below] at (2.2,-0.45){odlehlá hodnota};
    \end{scope}
    \begin{scope}[xshift=11.6cm]
        \draw[diredge] (0,0)--(4.4,0);
        \draw[diredge] (0,0)--(0,3.5);
        \draw[red!75!black,very thick] (0.5,0.8)--(4,2.5);
        \foreach \x/\y in {0.5/0.7,1/1.1,1.5/1.2,2/0.9,2/2.2,2.6/0.7,2.6/2.8,3.2/0.5,3.2/3.1,3.8/0.3,3.8/3.3}
            \fill[green!55!black](\x,\y) circle(2pt);
        \node[below,align=center] at (2.2,-0.45){proměnlivá\\velikost chyby};
    \end{scope}
\end{tikzpicture}

    \caption{Tři typické problémy lineárního modelu.}
    \label{fig:ai-regrese-omezeni}
\end{figure}

\begin{itemize}
    \item \textbf{Nelineární vztah.} Body mohou sledovat křivku, kterou žádná přímka nevystihne. Rezidua potom nevypadají náhodně, ale vytvářejí systematický obrazec.
    \item \textbf{Odlehlé hodnoty.} Metoda nejmenších čtverců zvýrazňuje velká rezidua. Jediný vzdálený bod proto může znatelně změnit směrnici i~absolutní člen.
    \item \textbf{Proměnlivá velikost chyby.} Rozptyl reziduí může s~rostoucím vstupem růst. Jediná hodnota \(\MSE\) pak zakrývá, že model je v~různých částech rozsahu různě přesný.
    \item \textbf{Závislá pozorování.} Hodnoty měřené v~čase nebo opakovaně u~stejného objektu mohou být vzájemně závislé. Běžná interpretace výsledků lineární regrese s~tím nemusí počítat.
    \item \textbf{Extrapolace.} Přímka je určena daty v~pozorovaném rozsahu. Daleko za tímto rozsahem se může skutečný vztah chovat zcela jinak.
    \item \textbf{Korelace není kauzalita.} Regresní přímka popisuje statistický vztah, ale sama neprokazuje, že změna \(\mathbf{X}\) způsobuje změnu \(\mathbf{Y}\).
\end{itemize}

\importantbox{Model není správný jen proto, že jej lze vypočítat. Výsledek je nutné posuzovat společně s~grafem dat, rezidui, způsobem sběru dat a~otázkou, pro kterou má být použit.}
